\documentclass[11pt,a4paper]{article}
\usepackage[margin=25mm]{geometry}
\usepackage{amsmath,amssymb,amsthm,hyperref}
\newtheorem{lemma}{Lemma}
\newtheorem{theorem}[lemma]{Theorem}
\newcommand{\gen}[1]{\langle#1\rangle}
\title{Centralizers of a primitive action on pairs}
\author{Human Compression project: a component of the L044 recovery}
\date{}
\begin{document}\maketitle
Let $\mathcal B$ be a partition of $2m$ points into $m$ pairs,
where $m>2$. Its full stabilizer is $V\rtimes S_m$, with
$V=\mathbb F_2^m$ acting by independent flips. We first prove
a statement about the actual permutation action, rather than
assuming that a commuting permutation induces a permutation of
$\mathcal B$.

\begin{lemma}\label{centralizer}
Suppose $B\le V\rtimes S_m$ is transitive on $\mathcal B$
and its induced group $P$ is primitive and generated by
involutions. Then either every element of $C_{S_{2m}}(B)$
preserves $\mathcal B$, or $B$ is regular on the $2m$ points,
$m$ is an odd prime, and $P\cong B$ is dihedral of order $2m$.
In the latter case an independent generating family of
involutions for $P$ has at most two members.
\end{lemma}
\begin{proof}
Put $N=B\cap V$. If $N\ne1$, some coordinate projection of
$N$ is nontrivial. Conjugation by $B$ permutes its coordinate
projections transitively, so every one is nontrivial. Its point
orbits are consequently exactly the pairs of $\mathcal B$.
A permutation centralizing $B$ centralizes each element of $N$
and therefore permutes these $N$-orbits. This proves the claim
in that case; no characteristic-subgroup assertion is needed.

Suppose $N=1$, and identify $B$ with $P$ via the block action.
A point orbit of $B$ meets every pair, and the size of that
intersection is constant. There are either two orbits, each
meeting every pair once, or one orbit of size $2m$.
In the first case each row action is the literal faithful
primitive action $P$ on $\mathcal B$. The centralizer of that
action is trivial. Indeed the orbits of a nontrivial centralizer
would form a $P$-invariant block system; primitivity would make
that centralizer transitive, forcing $P$ to be regular. A regular
primitive group is cyclic of prime order, and for $m>2$ is not
generated by involutions. Thus an equivariant permutation of
either row is the identity. An equivariant bijection between
the two rows is the unique bijection matching their points in
each original pair: compose any such bijection with that one
to obtain a centralizer of one row. Every permutation commuting
with $B$ therefore preserves the pairs.

In the transitive case fix a point $x$ in a pair $\Delta$.
Write $J=B_x$ and $M=B_\Delta$. The group $M$ is maximal in
$B$, because the block action is primitive, and $[M:J]=2$.
Hence $J\unlhd M$ and $M\le N_B(J)$. Maximality gives
$N_B(J)=M$ or $N_B(J)=B$.
For a faithful transitive action, the centralizing permutations
are the equivariant maps of the coset action on $B/J$; their
possible images of $J$ are precisely the cosets in $N_B(J)/J$.
If this normalizer is $M$, every such map sends $x$ within
$\Delta$, and equivariance sends every pair to itself.

If the normalizer is $B$, then $J\unlhd B$. Faithfulness of
the point action forces $J=1$, so $B$ is regular and $|M|=2$.
Let $a$ generate $M$. Some involution $b$ of $B$ lies outside
$M$, since the involutions generate $B$. Maximality of $M$
forces $B=\gen{a,b}$, a dihedral group. Its rotation subgroup
has order $m$. A proper nontrivial divisor of $m$ would give
a dihedral subgroup strictly between $M$ and $B$, contrary to
maximality. Thus $m$ is prime, and $m>2$ makes it odd.
Two distinct involutions generate this entire dihedral group,
so a family of at least three involutions cannot be independent.
\end{proof}

\begin{theorem}\label{proper}
Let $G=A_{2m}=\gen{\rho_0,\ldots,\rho_{r-1}}$ in its
natural action, with a string C-marking, and let $m\ge8$.
Suppose $H=G_i$ is transitive and preserves a maximal block
system consisting of $m$ pairs. If its primitive block image
$P$ is different from $A_m,S_m$, then $r\le m-1$.
Moreover, in every putative case $r\ge m$, regardless of $P$,
the connected block-generating interval $L$ used below has
at most one remaining noncommuting neighbour in $H$.
\end{theorem}
\begin{proof}
Choose an inclusion-minimal distinguished family $L$ whose
images generate $P$. Its images are independent and nonidentity
involutions. Write $C$ for the marks of $H$ outside $L$
commuting with all of $L$, and $R$ for the other marks.
We acquire the full L047 maximal-block decomposition.
If the noncommuting diagram of the images of $L$ is disconnected,
then $m\ge60$, $|L|\le2\log_2m$, $|C|\le1$ and $|R|\le4$.
It follows that
\[
 r=1+|L|+|C|+|R|\le2\log_2m+6<m.
\]
The last inequality holds at $m=60$ and its right-hand difference
is increasing for $m\ge60$.

Otherwise $L$ is an interval of the full string, $|C|\le1$,
and $|R|\le2$. Put $B=\gen L$ and assume $r\ge m$.
The exceptional alternative of Lemma~\ref{centralizer} would
give $|L|\le2$ and hence $r\le6<m$. Thus every permutation
centralizing $B$ preserves the chosen pairs. The omitted
$\rho_i$ cannot centralize $B$: together with $H$ it generates
$A_{2m}$, which preserves no nontrivial block system.
It therefore fails to commute with some mark of $L$.
By the string relations, $i$ is adjacent to that interval.
One of the only two possible noncommuting neighbours of $L$
is consequently absent from $H$, so $|R|\le1$ and
\[
                         r\le |L|+3.                 \tag{1}
\]
This proves the final assertion for every primitive block image.

For $P\ne A_m,S_m$, the accepted correction to the source's
Proposition~3.3 applies to the projected family: an independent
generating string of involutions in a primitive group of degree
$m\ge8$, other than $A_m,S_m$, has at most $m-4$ members.
Its statement does not require the intersection property in the
quotient. Thus $|L|\le m-4$, and (1) gives $r\le m-1$,
a contradiction to the supposed exception. This proves the theorem.
\end{proof}

\paragraph{Inputs and scope.}
The source is Cameron--Fernandes--Leemans--Mixer,
\emph{Highest rank of a polytope for $A_n$},
arXiv:1605.09173v1, Section~3.2. The L047 block decomposition
and its precise classical inputs remain acquired. The projected
string bound is the accepted corrected theorem, with its full
finite proof and declared software premises; the source's false
bound for arbitrary independent families is not used.
Lemma~\ref{centralizer} supplies the missing justification for
attaching the omitted label to $L$ in this pair-block branch.
It uses the actual block kernel and the normalizer of a point
stabilizer; it does not treat a block-breaking permutation as
already acting on the blocks. The cases $P=A_m,S_m$ remain
separate, and the degree-twelve and fourteen boundaries are
covered by acquired complete finite classifications. No new
computation, whole-original credit, conjecture answer or Lean
verification is asserted here.
\end{document}
